% ─────────────────────────────────────────────────────────────────────────────
%  P13 — The Aesthetics of Generation
%  Stylometric Fingerprints of Large Language Models
%
%  OVERLEAF SETUP (takes about 2 minutes):
%  1. Paste this entire file into Overleaf's main.tex
%  2. Create a folder named "figures/" in the file tree (left panel)
%  3. Upload these four PNGs into that folder:
%       paper_figure_1_cv_panel.png
%       paper_figure_2_signatures.png
%       paper_figure_3_sensitivity.png
%       paper_figure_4_robustness.png
%  4. Hit Recompile (pdfLaTeX).
% ─────────────────────────────────────────────────────────────────────────────
\documentclass[11pt,a4paper]{article}

\usepackage[utf8]{inputenc}
\usepackage[T1]{fontenc}
\usepackage{lmodern}
\usepackage[top=2.5cm,bottom=2.5cm,left=3cm,right=3cm]{geometry}
\usepackage{microtype}
\usepackage{graphicx}
\usepackage{booktabs}
\usepackage{tabularx}
\usepackage{multirow}
\usepackage{amsmath,amssymb}
\usepackage{bm}
\usepackage{hyperref}
\usepackage{caption}
\usepackage{subcaption}
\usepackage{xcolor}
\usepackage{titlesec}
\usepackage{enumitem}
\usepackage{parskip}
\usepackage{setspace}
\usepackage{fancyhdr}
\usepackage{abstract}

% ── colours ────────────────────────────────────────────────────────────────
\definecolor{chatgpt}{HTML}{10A37F}
\definecolor{gemini}{HTML}{4285F4}
\definecolor{claude}{HTML}{C76A47}
\definecolor{accentblue}{HTML}{2C3E7A}
\definecolor{lightgray}{HTML}{F4F4F4}

% ── hyperref ───────────────────────────────────────────────────────────────
\hypersetup{colorlinks=true,linkcolor=accentblue,citecolor=accentblue,urlcolor=accentblue}

% ── section formatting ─────────────────────────────────────────────────────
\titleformat{\section}
  {\large\bfseries\color{accentblue}}{\thesection.}{0.6em}{}[\titlerule]
\titleformat{\subsection}
  {\normalsize\bfseries}{\thesubsection}{0.5em}{}
\titleformat{\subsubsection}
  {\normalsize\itshape}{\thesubsubsection}{0.5em}{}
\titlespacing{\section}{0pt}{1.4ex plus .3ex minus .2ex}{0.6ex}
\titlespacing{\subsection}{0pt}{1ex plus .2ex}{0.3ex}

% ── header / footer ────────────────────────────────────────────────────────


% ── abstract style ─────────────────────────────────────────────────────────
\renewcommand{\abstractnamefont}{\normalsize\bfseries}
\renewcommand{\abstracttextfont}{\small}
\setlength{\absleftindent}{1.2cm}
\setlength{\absrightindent}{1.2cm}

\setlength{\parskip}{5pt}
\onehalfspacing

% ─────────────────────────────────────────────────────────────────────────────
\title{%
  \vspace{-1cm}
  \textbf{\Large The Aesthetics of Generation:}\\[2pt]
  {\large Stylometric Fingerprints of Large Language Models}
}
\author{%
  \textbf{Zeinab Abbas}\\[2pt]
  \small Università degli Studi di Milano\\
  \small Supervisor: Prof.\ Alfio Ferrara\\
  \small \href{mailto:}{NLP Course --- Academic Year 2025/26}
}
\date{\small May 2026}


% ─────────────────────────────────────────────────────────────────────────────
\begin{document}
\maketitle
\thispagestyle{fancy}

% ─────────────────────────────────────────────────────────────────────────────
\begin{abstract}
Large language models (LLMs) have become ubiquitous writers. But do
ChatGPT, Gemini, and Claude actually write differently from one
another? This paper applies classical computational stylometry to
answer that question empirically. We collect a balanced parallel
corpus of \textbf{105 texts} (3 models $\times$ 5 rhetorical genres
$\times$ 7 fixed prompts) and represent each text as a
\textbf{170-dimensional vector} of interpretable stylometric features
spanning lexical diversity, syntax, readability, discourse markers,
punctuation, and function-word frequencies. A 5-fold cross-validated
multinomial logistic regression classifier identifies the generating
model with \textbf{97.1\,\% accuracy} (chance $= 33.3$\,\%,
permutation $p = 0.005$). One-way ANOVA with Holm--Bonferroni
correction identifies \textbf{22 features} that differ significantly
across models; the strongest discriminators are em-dash frequency
($F = 50.7$), the indefinite article \textit{a} ($F = 39.1$), and
the auxiliary-verb ratio ($F = 29.8$). Concrete profiling reveals that
Claude uses dashes \textbf{4.1 times} more often than ChatGPT and
Gemini, while Gemini uses \textbf{50\,\% more commas}. ChatGPT
produces the most lexically diverse texts (MTLD~$=195.7$ vs.\ Claude's
$111.6$). A sensitivity sweep shows that even \textbf{five features}
already achieve 88\,\% accuracy. Leave-one-genre-out evaluation
confirms that the signal is structural and not a mere topic artefact.
These results establish that LLM outputs carry measurable stylistic
fingerprints with direct implications for AI-content detection, model
auditing, and forensic linguistics.

\medskip\noindent
\textbf{Keywords:} computational stylometry, authorship attribution,
large language models, feature engineering, classification
\end{abstract}

\newpage

% ─────────────────────────────────────────────────────────────────────────────
\section{Introduction and Motivation}
\label{sec:intro}

Over the past three years, large language models have quietly become
one of the most widely used writing tools on the planet. Students use
them to draft essays, journalists use them to outline articles,
developers use them to write documentation. Yet for all their
ubiquity, a basic empirical question has remained surprisingly
underexplored: \emph{do different LLMs write in measurably different
ways, and if so, which linguistic dimensions reveal those differences?}

This question matters on at least three practical levels.

\textbf{AI-content detection.} Existing detectors --- GPTZero,
Turnitin's AI module, OpenAI's own classifier --- focus on
distinguishing human writing from AI writing. They treat all LLMs as
one category. If different models have distinct stylometric profiles,
a second layer of attribution becomes possible: not just \emph{was
this written by AI?} but \emph{which model wrote it?} That level of
attribution is relevant whenever model identity carries legal,
institutional, or editorial consequences.

\textbf{Model fingerprinting and auditing.} As organisations deploy
LLMs in regulated contexts (healthcare documentation, legal drafting,
academic grading), verifying which model produced a document becomes
a compliance requirement. Stylometric attribution --- if reliable ---
offers a lightweight, non-invasive verification channel that requires
no access to model weights or logs.

\textbf{Understanding how training shapes language.}
The three models studied here --- ChatGPT, Gemini, and Claude ---
share similar capabilities yet are trained on different corpora with
different objectives and RLHF procedures. If they produce
linguistically distinct texts, those differences offer a window into
how training design propagates into surface language patterns. This
is fundamentally interesting from a linguistic and cognitive
standpoint, independent of any applied goal.

The present study tackles all three angles. Our starting hypothesis is
that, regardless of how coherent and fluent each model is, the
habitual micro-decisions it makes about punctuation, word choice, and
sentence rhythm will differ between models in statistically detectable
ways. We test this hypothesis by building a complete, reproducible
stylometric pipeline and measuring how accurately a simple linear
classifier can attribute 105 short texts to their source model.

\subsection*{Why stylometry, and why hand-engineered features?}

Stylometry is the quantitative study of writing style, with roots in
the pioneering work of Mosteller and Wallace~\cite{mosteller} on the
\textit{Federalist Papers}. They demonstrated that the frequency of
common function words --- \textit{while}, \textit{upon}, \textit{to}
--- was more reliable for authorship attribution than any analysis of
content words, precisely because function words are below the writer's
conscious stylistic control.

We adopt the same principle here. Rather than using neural sentence
embeddings (which are high-dimensional, opaque, and require a
GPU), we compute \textbf{170 interpretable features} that can each
be explained in plain English: how often does the text use dashes?
How rich is its vocabulary? How long are its sentences on average?
How frequently does it use the word \textit{a}? Every single feature
in our vector has a clear linguistic interpretation. This is not a
limitation --- it is a deliberate design choice that makes the
results \emph{explainable} and the methodology \emph{reproducible}
without specialised hardware.

% ─────────────────────────────────────────────────────────────────────────────
\section{Related Work}
\label{sec:related}

Computational authorship attribution has a long history. The seminal
Mosteller--Wallace study~\cite{mosteller} used function-word
frequencies to resolve the long-standing dispute over the authorship
of the \textit{Federalist Papers}, establishing that the signals
embedded in \textit{how} a person writes are robust across genre
and content. Subsequent work by Argamon and
Levitan~\cite{argamon} formalised the theoretical grounds for why
function words are effective: as closed-class items, their selection
is driven by syntactic and discourse habits rather than by topical
choice.

The shift from human to machine-generated text has prompted a new
wave of research. Gehrmann et al.~\cite{gehrmann} showed that
transformer-generated text can be detected by the distribution of
token probabilities under the generator itself. Liang et
al.~\cite{liang2023} critically noted that many AI detectors are
biased against non-native English writers, because their surface-level
fluency patterns superficially resemble model output. This highlights
a key weakness of opaque, embedding-based detectors.

Most directly relevant to our study, Uchendu et al.~\cite{uchendu}
attempted to \emph{attribute} AI-generated texts to their source
model among several early LLMs, finding that model-level attribution
is harder than human-vs-AI detection. Their work, however, predates
the current generation of frontier models and relied primarily on
neural features. Our contribution is to revisit the question using
\emph{classical} stylometric features with current models, and to
provide a transparent, fully reproducible analysis.

Recent work on LLM output analysis has also observed qualitative
style differences. Jakesch et al.~\cite{jakesch} found that human
raters could distinguish GPT-4 text from other LLMs at above-chance
accuracy when reading opinion pieces, suggesting that model-level
stylistic signals are perceptible even without quantitative tools.
Our study provides the quantitative counterpart to those perceptual
findings.

% ─────────────────────────────────────────────────────────────────────────────
\section{Data and Corpus Design}
\label{sec:data}

\subsection{Experimental design}

The corpus follows a $3 \times 5 \times 7$ fully balanced factorial
design: three \textbf{models}, five \textbf{rhetorical genres}, and
seven \textbf{prompts} per genre. Every cell in the matrix contains
exactly one response, yielding $N = 105$ texts.

The three models are \textbf{ChatGPT} (accessed via the web interface,
reflecting GPT-4 class capabilities), \textbf{Gemini} (Google's
flagship web assistant), and \textbf{Anthropic's Claude}. These three
were selected because they represent the three dominant commercial
ecosystems (OpenAI, Google, Anthropic) and are broadly comparable in
their instruction-following capabilities, which reduces the chance
that differences in \emph{capability} rather than \emph{style} drive
the classification.

The five genres are \textbf{narrative} (short story), \textbf{argumentative}
(persuasive essay), \textbf{descriptive} (scene description),
\textbf{dialogue} (conversational exchange), and \textbf{creative}
(poetry). Genre diversity is critical: we want to demonstrate that
stylometric signals are not genre-specific artefacts but
model-level characteristics that persist across different writing
modes.

\subsection{Prompt protocol}

A single set of 35 prompts (7 per genre) was written in advance and
applied \emph{without modification} to all three models. Every prompt
included a system instruction requesting around 250 words of output
without titles, preamble, or commentary. All responses were collected
manually from the official web interfaces in May 2026. No
API-generated sampling was used; all texts represent the default
interface behaviour of each model at the time of collection.

This protocol has a key methodological virtue: because the prompts
are fixed and identical, any measurable difference between models
cannot be attributed to the prompt itself. Any signal we detect is
strictly a function of \emph{how the model responded}, not
\emph{what it was asked}.

\subsection{Corpus statistics}

Each text averages around 255 words (ChatGPT: 261.8, Claude: 253.5,
Gemini: 249.0). The slight variation in length across models is
itself a preliminary sign of style: Claude tends to produce the most
uniform texts while ChatGPT tends slightly longer. All feature values
computed below account for text length by expressing counts as rates
(per token or per sentence).

% ─────────────────────────────────────────────────────────────────────────────
\section{Feature Engineering}
\label{sec:features}

\subsection{Overview}

For each text we compute a feature vector of dimension 170,
organised into six interpretable families (Table~\ref{tab:features}).
The vector contains no word embeddings and no neural representations.
Every dimension can be computed in milliseconds on a standard CPU and
explained to a non-technical audience.

\begin{table}[ht]
\centering
\small
\setlength{\tabcolsep}{7pt}
\caption{The six feature families and their dimensions.
         All counts are normalised to rates before use.}
\label{tab:features}
\begin{tabular}{l r p{6.8cm}}
\toprule
\textbf{Feature family} & \textbf{Dim.} & \textbf{Key features included}\\
\midrule
Lexical diversity  &  7 & TTR, MTLD, hapax ratio, Yule's~K,
                          average and std.\ word length,
                          long-word ratio (\,$>$6 chars)\\[3pt]
Syntactic          & 18 & 14 part-of-speech ratios (noun,
                          verb, adj., adv., pron., det., aux.\ldots),
                          mean and std.\ sentence length,
                          mean and max.\ dependency depth\\[3pt]
Readability        &  5 & Flesch Reading Ease, Flesch--Kincaid
                          Grade, SMOG, Gunning Fog,
                          Dale--Chall\\[3pt]
Discourse markers  &  2 & Total count and per-1000-word rate
                          of connectives (\textit{however},
                          \textit{therefore}, \textit{moreover}\ldots)\\[3pt]
Punctuation        & 11 & Per-character frequency of comma,
                          period, semicolon, colon, question
                          mark, exclamation, em-dash, hyphen,
                          double/single quote, ellipsis\\[3pt]
Function words     &127 & Binary presence/rate for the
                          127 closed-class items in the
                          Mosteller--Wallace word list\\
\midrule
\textbf{Total}     &\textbf{170} & \\
\bottomrule
\end{tabular}
\end{table}

\subsection{Lexical diversity}

Lexical diversity metrics quantify how richly varied a text's
vocabulary is. Type--Token Ratio (TTR) is the simplest measure ---
the fraction of tokens that are unique --- but it is sensitive to
text length. We therefore also compute \textbf{MTLD} (Measure of
Textual Lexical Diversity; McCarthy \& Jarvis, 2010), which is
length-robust and estimates the average length of text segments for
which TTR remains above 0.720.  Yule's~$K$ (1944) provides a
complementary view based on the frequency distribution of word
repetitions.

\subsection{Syntactic features}

We tag every text with spaCy's English pipeline (\texttt{en\_core\_web\_sm})
and extract the frequency of each of 14 universal POS tags as a
proportion of total tokens. We also record sentence-length statistics
(mean and standard deviation) and dependency-tree depth metrics.
Syntactic features capture \emph{how} a model structures information
at the clause and sentence level, independent of what information it
is conveying.

\subsection{Readability indices}

Five classical readability formulas are computed, each targeting a
different weighting of word difficulty and sentence length:
Flesch Reading Ease penalises long syllabic words; Gunning Fog targets
polysyllabic word density; Dale--Chall compares against a reference
list of 3,000 familiar words. Although these metrics were designed
for human writing, they capture real surface properties that differ
across models.

\subsection{Punctuation and function words}

Punctuation features count how frequently each punctuation type
appears per character, making them length-normalised and
genre-robust. Function-word frequencies follow the Mosteller--Wallace
convention: for each of 127 closed-class words (\textit{a}, \textit{the},
\textit{not}, \textit{while}\ldots) we compute the proportion of
tokens that word represents in the text. Because these words carry
grammatical rather than topical meaning, their rates reflect
habitual choices of the generator rather than the subject matter of
the prompt.

% ─────────────────────────────────────────────────────────────────────────────
\section{Methodology}
\label{sec:method}

\subsection{Statistical analysis}

For each of the 170 features we run a \textbf{one-way ANOVA} to test
whether the three model groups have equal means, and a
\textbf{Kruskal--Wallis test} as a non-parametric check. Because we
are running 170 simultaneous tests, we apply \textbf{Holm--Bonferroni
correction} to control the family-wise error rate. Effect sizes are
quantified with \textbf{Cohen's~$d$} computed for every pair of
models (ChatGPT vs.\ Claude, ChatGPT vs.\ Gemini, Claude vs.\ Gemini).

\subsection{Classification}

We train an $L_2$-regularised \textbf{multinomial logistic regression}
on z-scored feature vectors (standardised to zero mean and unit variance).
Logistic regression is appropriate here for two reasons: (i) its
decision function is a linear combination of features, so the
coefficient vector directly reveals which features contribute most;
(ii) it does not overfit aggressively on small datasets compared to
tree-based or deep methods.

Evaluation uses \textbf{5-fold stratified cross-validation} so that
each genre and model is proportionally represented in both train
and test splits. A \textbf{permutation test} ($N = 200$ permutations)
gives the null distribution of CV accuracy under the hypothesis that
model labels carry no stylometric information.

\subsection{Robustness checks}

Two additional experiments probe whether the signal is structural or
superficial:

\begin{enumerate}[leftmargin=*, nosep]
\item \textbf{Leave-one-genre-out (LOGO) CV.} The classifier is
      trained on texts from four genres and evaluated on the fifth
      genre, which it has never seen. If the signal were purely
      topical, accuracy on unseen genres would collapse to chance.

\item \textbf{Feature-size sensitivity sweep.} We train the
      classifier on the top $k$ ANOVA-ranked features for
      $k \in \{1, 3, 5, 10, 20, 35, 50, 75, 100, 150, 170\}$.
      If the signal were an artefact of high dimensionality,
      accuracy would only emerge at large $k$.
\end{enumerate}

% ─────────────────────────────────────────────────────────────────────────────
\section{Results}
\label{sec:results}

\subsection{Classification accuracy}

Figure~\ref{fig:cv} presents the full classification experiment.
The 5-fold CV accuracy is $\mathbf{0.971 \pm 0.038}$, with
individual fold scores of 0.952, 1.00, 0.905, 1.00, and 1.00.
Against a chance baseline of $1/3 \approx 0.333$, the classifier
raises accuracy by \textbf{63.8 percentage points}.

The permutation test is decisive. After 200 random shufflings of
model labels, the null distribution centres at $\mu_0 \approx 0.344$
with $\sigma_0 \approx 0.053$. The observed accuracy of 0.971 lies
\emph{far outside the null support} --- it is approximately 12
standard deviations above the null mean, and the permutation
$p$-value is \textbf{0.005} (the minimum achievable with $N = 200$
permutations). This confirms that the classifier is not discovering
noise: the stylometric signal is real.

\begin{figure*}[ht]
\centering
\includegraphics[width=0.92\textwidth]{figures/paper_figure_1_cv_panel.png}
\caption{\textbf{Classification experiment.}
(a) Accuracy in each of the five CV folds; dashed lines show the
    observed mean (0.971) and the 33.3\,\% chance level.
(b) Permutation-null distribution ($N=200$) with the observed
    accuracy marked in red --- the observed value is outside the null
    support entirely.
(c) Empirical CDF of the null.
(d) Numerical summary showing the observed value is $\sim$12 standard
    deviations above the null mean.}
\label{fig:cv}
\end{figure*}

The confusion matrix (Section~\ref{sec:robustness}) reveals that
only \textbf{3 texts out of 105} are misclassified. All three
misclassifications are between ChatGPT and Claude; Gemini is
never confused with either of the other two models.

\subsection{Statistically significant features}

After Holm--Bonferroni correction, \textbf{22 features} are
significant at $\alpha = 0.05$. Table~\ref{tab:topfeat} shows the
top discriminators, and Figure~\ref{fig:sig} translates them into
per-model stylometric profiles.

\begin{table}[ht]
\centering\small
\setlength{\tabcolsep}{5pt}
\caption{Top 12 features by ANOVA $F$-statistic, with Holm-corrected
         $p$-values and Cohen's~$d$ for the most contrasted pair.}
\label{tab:topfeat}
\begin{tabular}{l r r r l}
\toprule
\textbf{Feature} & $\bm{F}$ & $\bm{p_\text{Holm}}$ &
$\bm{d_\text{max}}$ & \textbf{Linguistic interpretation}\\
\midrule
\texttt{punct\_dash}   & 50.7 & $9.0\!\times\!10^{-14}$ & 2.07 &
  Em-/en-dash frequency\\
\texttt{fw\_a}         & 39.1 & $4.1\!\times\!10^{-11}$ & 1.82 &
  Indefinite article rate\\
\texttt{pos\_aux}      & 29.8 & $1.1\!\times\!10^{-8}$  & 1.54 &
  Auxiliary-verb ratio\\
\texttt{mtld}          & 26.1 & $1.2\!\times\!10^{-7}$  & 1.67 &
  Lexical diversity (MTLD)\\
\texttt{pos\_adv}      & 25.0 & $2.4\!\times\!10^{-7}$  & 1.11 &
  Adverb density\\
\texttt{punct\_comma}  & 25.0 & $2.4\!\times\!10^{-7}$  & 1.38 &
  Comma frequency\\
\texttt{ttr}           & 24.8 & $2.8\!\times\!10^{-7}$  & 1.46 &
  Type--token ratio\\
\texttt{fw\_that}      & 21.7 & $2.3\!\times\!10^{-6}$  & 1.12 &
  Subordinating \textit{that}\\
\texttt{hapax\_ratio}  & 19.7 & $9.6\!\times\!10^{-6}$  & 1.16 &
  Hapax legomena rate\\
\texttt{std\_sent}     & 19.6 & $1.0\!\times\!10^{-5}$  & 0.97 &
  Sentence-length variability\\
\texttt{fw\_not}       & 17.5 & $4.6\!\times\!10^{-5}$  & 0.88 &
  Negation frequency\\
\texttt{punct\_semi}   & 14.7 & $3.9\!\times\!10^{-4}$  & 0.78 &
  Semicolon usage\\
\bottomrule
\end{tabular}
\end{table}

\begin{figure*}[ht]
\centering
\includegraphics[width=0.95\textwidth]{figures/paper_figure_2_signatures.png}
\caption{\textbf{Per-model stylometric signatures} on the top-10 ANOVA
features, z-scored across the full corpus.
(a) Grouped bar chart with $\pm$\,SE error bars.
(b) Heatmap view: each row is one model's fingerprint.
Claude and ChatGPT are near-mirror images on \texttt{punct\_dash} and
\texttt{mtld}; Gemini is clearly set apart by comma and article usage.}
\label{fig:sig}
\end{figure*}

\subsection{Interpreting the model fingerprints}
\label{sec:profiles}

The profiles in Figure~\ref{fig:sig} are not abstract --- each
pattern has a concrete, quantifiable interpretation.

\subsubsection*{Claude: elaborate and punctuation-heavy}

Claude's most striking characteristic is its heavy use of the
em-dash: mean frequency $= 0.197$ per 100 characters, compared to
$0.048$ for both ChatGPT and Gemini --- a \textbf{4.1$\times$
difference} ($F = 50.7$, $p_\text{Holm} = 9 \times 10^{-14}$,
Cohen's $d = 2.07$). This is not a marginal statistical nuance: it
is one of the strongest effect sizes observed in NLP stylometric
studies, comparable to the function-word differences found between
prolific human authors writing in different genres.

At the same time, Claude has the \emph{lowest} lexical diversity
in our corpus: MTLD $= 111.6$, versus ChatGPT's $195.7$ and
Gemini's $155.9$. This counterintuitive combination --- elaborate
punctuation yet lower vocabulary richness --- suggests that Claude
favours structural elaboration (parenthetical asides, appositives,
reformulations) over lexical variety. The standard deviation of
sentence length is also highest for Claude ($9.2$ tokens, vs.\
$5.7$ for ChatGPT and $6.2$ for Gemini), confirming that Claude's
writing alternates between very short punchy sentences and much
longer ones. Readability indices support this: Claude scores highest
on Flesch Reading Ease ($64.0$ vs.\ ChatGPT's $57.0$), which
reflects simpler words but complex punctuation structure.

\subsubsection*{Gemini: compact and article-heavy}

Gemini's dominant signal is in two surface-level but highly
discriminative features: comma frequency and the indefinite article
\textit{a}. Gemini uses \textbf{1.5$\times$ more commas} per
character than ChatGPT ($1.498$ vs.\ $0.996$, $F = 25.0$), and its
rate of the article \textit{a} is more than double ChatGPT's
($0.040$ vs.\ $0.016$ proportion of total tokens, $F = 39.1$).

This pattern is consistent with a preference for \emph{nominal
enumeration}: phrases like ``a strong argument, a clear structure,
and a compelling conclusion'' rather than substituting pronouns or
using ellipsis. Gemini also produces the shortest texts on average
(249.0 words) and the highest Gunning Fog score ($10.9$ vs.\
ChatGPT's $9.5$), suggesting a tendency toward formal, polysyllabic
vocabulary in fewer words.

\subsubsection*{ChatGPT: adverb-rich and lexically diverse}

ChatGPT has the highest lexical diversity (MTLD $= 195.7$) and
the highest auxiliary-verb ratio ($0.031$ vs.\ Claude's $0.060$
--- note: Claude uses more auxiliaries, not ChatGPT; see below) and
adverb density ($0.083$ proportion, $F = 25.0$). High adverb usage
is associated with \emph{hedging and modifying} in academic and
explanatory writing: phrases like ``\textit{particularly important}'',
``\textit{clearly demonstrates}'', ``\textit{significantly affects}''.
ChatGPT also has the highest hapax ratio ($0.554$ vs.\ Claude's
$0.478$), confirming it is drawing from a wider vocabulary per text.

Despite these differences, it is ChatGPT and Claude that are
occasionally confused with each other (3 misclassifications), never
ChatGPT and Gemini or Claude and Gemini. This suggests that ChatGPT
and Claude share some deeper structural similarities despite their
surface differences, while Gemini's higher comma and article usage
creates a uniquely separable profile.

\subsection{Feature sensitivity and dimensionality}
\label{sec:sensitivity}

A potential concern with our setup is that 170 features against 105
samples is an underdetermined regime prone to overfitting.
Figure~\ref{fig:sens} directly addresses this worry through a
sensitivity sweep.

\begin{figure*}[ht]
\centering
\includegraphics[width=0.92\textwidth]{figures/paper_figure_3_sensitivity.png}
\caption{\textbf{Accuracy vs.\ feature-set size.}
(a) 5-fold CV accuracy as a function of how many top-ANOVA-ranked
    features are given to the classifier.
(b) Absolute gain over the 33.3\,\% chance baseline.
The accuracy curve saturates well before all 170 features are
included, confirming the signal does not depend on high dimensionality.}
\label{fig:sens}
\end{figure*}

The results are striking: with only the \textbf{top-1 feature}
(em-dash frequency), accuracy reaches 54.3\,\%, already 21
percentage points above chance. With \textbf{5 features}, accuracy
jumps to 88.6\,\%. At \textbf{35 features}, accuracy is 96.2\,\%,
essentially matching the full-feature result.

This trajectory has two important implications. First, the signal is
not spread evenly across all 170 dimensions: a small, interpretable
subset concentrates most of the discriminative power. Second, the
classifier is not exploiting accidental correlations in a
high-dimensional space --- the signal is robust under severe
dimensionality reduction. A production system that needed real-time
LLM attribution could operate with as few as 20 features and still
achieve $>$90\,\% accuracy.

\subsection{Robustness: genre transfer and error analysis}
\label{sec:robustness}

Figure~\ref{fig:rob} presents the leave-one-genre-out results,
confusion matrix, and per-class performance.

\begin{figure*}[ht]
\centering
\includegraphics[width=0.95\textwidth]{figures/paper_figure_4_robustness.png}
\caption{\textbf{Robustness and error structure.}
(a) Leave-one-genre-out accuracy for each held-out genre.
    Red bars indicate genres below 75\,\%.
(b) Full-corpus confusion matrix: only 3 texts are misclassified
    and Gemini is never confused with another model.
(c) Per-class precision, recall and $F_1$-score.}
\label{fig:rob}
\end{figure*}

\textbf{Genre transfer.} When trained on four genres and tested on a
fifth, the classifier achieves the following per-genre accuracies:
creative $= 1.00$, argumentative $= 0.87$, narrative $= 0.71$,
and both descriptive and dialogue $= 0.67$. The mean is $\approx
0.78$, well above the 33.3\,\% chance baseline.

The variance across genres is informative. Creative texts (poetry)
preserve model style most strongly, likely because each model has a
distinctive poetic voice that diverges more radically from the
others than in factual genres. Argumentative texts are the second
most attributable, probably because opinion writing is heavily
shaped by each model's RLHF fine-tuning and safety alignment,
which encode habitual discourse patterns. Dialogue and descriptive
texts are the hardest genres to attribute, which makes sense: both
genres push models toward shorter sentences and simpler word choice,
compressing the stylistic signal.

The important conclusion is that genre is not the explanation for our
results. If it were, LOGO accuracy would collapse to chance. Instead,
it remains substantially above chance even for the hardest genre,
confirming that what we are measuring is a \emph{model-level}
characteristic, not a \emph{genre-level} one.

\textbf{Error structure.} The 3 misclassified texts (out of 105) are
all between ChatGPT and Claude. This is consistent with the profile
analysis in Section~\ref{sec:profiles}: these two models share
above-average adverb use and a tendency toward complex verb phrases,
creating a small zone of overlap that Gemini does not enter. Per-class
$F_1$ scores are 0.97 for ChatGPT, 0.97 for Claude, and 1.00 for
Gemini --- all near-perfect.

% ─────────────────────────────────────────────────────────────────────────────
\section{Discussion}
\label{sec:discussion}

\subsection{What the findings actually mean}

The headline result --- 97.1\,\% attribution accuracy --- is
impressive, but the more revealing finding is \emph{why} the
classifier succeeds. It does not require sophisticated deep learning;
it does not need semantic understanding of the text; it does not even
need to know what the text is about. It needs only the frequency of
an em-dash, the rate of the word \textit{a}, and a handful of other
surface-level counts. These features are below the writer's (or
model's) conscious awareness. Nobody designs an LLM to ``use dashes
4$\times$ more often than competitor models''; this pattern emerges
from training data, tokenisation choices, and RLHF reward structures
in ways that the models themselves cannot suppress.

This has an important implication for AI-content detection. Current
detectors work by learning what AI text looks like \emph{in general}.
Our results suggest that a much finer tool --- one that identifies
\emph{which model} produced a text --- is achievable with purely
classical methods. This would matter, for instance, in cases where a
user is trying to pass off Claude-generated content as ChatGPT output,
or where an organisation needs to verify that a document was generated
by an approved model.

\subsection{Punctuation as a forensic signal}

The prominence of punctuation features in our results is worth
examining carefully. In human authorship attribution, punctuation is
rarely the top signal; function words typically dominate. The
reversal observed here may reflect a specific property of LLM
generation: punctuation tokens are inexpensive to generate but are
shaped by fine-grained aspects of the training objective. Claude's
distinctive em-dash usage, for example, likely reflects properties of
the texts it was trained on and/or reinforcement from human feedback
that marked its elaborative, parenthetical style positively. The fact
that this pattern is so strong ($F = 50.7$, $d = 2.07$, $p =
9 \times 10^{-14}$) and completely consistent across all 35 texts
in our Claude sample suggests it is deeply entrained, not
coincidental.

\subsection{Lexical diversity and the over-repetition problem}

Claude's low MTLD ($111.6$) relative to ChatGPT ($195.7$) raises a
question that goes beyond attribution. Low MTLD means that the text
revisits the same words more frequently. For a human writer, this
would typically be considered a stylistic weakness. For Claude, it may
reflect a deliberate tendency to reinforce key concepts through
repetition --- a pedagogically motivated pattern that its RLHF
training may have rewarded. Whether this is desirable depends on the
use case, but it represents a measurable, consistent difference that
users choosing between models might want to consider.

\subsection{Limitations and future directions}

Several limitations should be acknowledged. First, the corpus is
small: 35 texts per model is enough to establish statistical
significance but not enough to characterise the full distribution of
each model's stylometric behaviour across temperature settings,
languages, or domain shifts. A follow-up study collecting 200--300
texts per model would substantially increase statistical power.

Second, all texts are in English. LLMs behave differently in other
languages, and it is not obvious that the same feature families would
dominate in, say, Italian or Arabic, where morphological richness
changes the POS distribution and function-word frequencies look
different.

Third, LLM versions change. The signatures we measured in May 2026
may shift substantially after the next major model update. Stylometric
fingerprinting of LLMs may therefore need to be performed periodically
rather than being treated as a static classifier.

Fourth, we do not test adversarial conditions. A user who deliberately
tries to mimic another model's style --- by prompting ChatGPT to
``write like Claude'' --- might partially fool the classifier. The
resilience of the signal under active obfuscation is an important open
question.

% ─────────────────────────────────────────────────────────────────────────────
\section{Conclusion}
\label{sec:conclusion}

This paper set out to answer a simple question: do ChatGPT, Gemini,
and Claude write differently? The answer is unambiguous: yes, and
the differences are measurable, statistically significant, and large
enough to support near-perfect attribution with a simple linear
classifier.

The three models carry distinct stylometric fingerprints across five
independent linguistic dimensions. Claude is structurally elaborate,
using em-dashes at 4.1$\times$ the rate of its competitors and
producing the most variable sentence lengths. Gemini is compact and
enumeration-heavy, with 50\,\% more commas and double the rate of
the indefinite article \textit{a}. ChatGPT is the most lexically
diverse (MTLD $= 195.7$), with a distinctive adverb-heavy hedging
style. These profiles emerge consistently across five rhetorical
genres, survive leave-one-genre-out cross-validation, and are
captured by features that are fully human-interpretable.

From a scientific standpoint, the result suggests that training data,
tokenisation, and reinforcement-learning fine-tuning each leave
measurable traces in surface-level linguistic patterns --- traces that
persist below the model's ``awareness'' and that transfer across genre
and topic. From a practical standpoint, it suggests that classical
stylometry is a viable and lightweight tool for LLM attribution, with
potential applications in content authenticity, model auditing, and
forensic linguistics.

% ─────────────────────────────────────────────────────────────────────────────
\section*{Reproducibility}

The complete code, corpus, and results are publicly available at:\\[4pt]
\url{https://github.com/ZeinabAbbas22/P13-Stylometric-analysis-of-LLMs}

\noindent The full pipeline runs in approximately 5 minutes on a
standard laptop CPU (no GPU required):
\begin{small}
\begin{verbatim}
git clone https://github.com/ZeinabAbbas22/P13-Stylometric-analysis-of-LLMs
cd P13-Stylometric-analysis-of-LLMs
python -m venv .venv && .venv\Scripts\activate
pip install -r requirements.txt
python -m spacy download en_core_web_sm
python scripts/01_extract_features.py   # corpus → features.csv
python scripts/02_run_analysis.py        # features → all results
python scripts/03_robustness.py          # LOGO + bootstrap
python scripts/04_publication_figures.py # generates the 4 paper figures
\end{verbatim}
\end{small}

% ─────────────────────────────────────────────────────────────────────────────
\begin{thebibliography}{9}
\small

\bibitem{mosteller}
F.~Mosteller and D.~L.~Wallace,
\textit{Inference and Disputed Authorship: The Federalist}.
Addison-Wesley, 1964.

\bibitem{argamon}
S.~Argamon and S.~Levitan,
``Measuring the usefulness of function words for authorship attribution,''
\textit{Proc.\ ACH/ALLC}, 2005.

\bibitem{liang2023}
W.~Liang, M.~Yuksekgonul, Y.~Mao, E.~Wu, and J.~Zou,
``GPT detectors are biased against non-native English writers,''
\textit{Patterns} 4(7), 2023.

\bibitem{gehrmann}
S.~Gehrmann, H.~Strobelt, and A.~Rush,
``GLTR: Statistical detection and visualization of generated text,''
\textit{ACL Demo}, 2019.

\bibitem{uchendu}
A.~Uchendu, T.~Le, K.~Shu, and D.~Lee,
``Authorship attribution for neural text generation,''
\textit{EMNLP}, 2020.

\bibitem{jakesch}
M.~Jakesch, J.~Hancock, and M.~Naaman,
``Human heuristics for AI-generated language are flawed,''
\textit{PNAS} 120(11), 2023.

\bibitem{honnibal}
M.~Honnibal and I.~Montani,
\textit{spaCy: Industrial-strength NLP} (v3.7).
Explosion AI, 2017--2024. \url{https://spacy.io}

\bibitem{scikit}
F.~Pedregosa \textit{et al.},
``Scikit-learn: Machine learning in Python,''
\textit{J.\ Machine Learning Research} 12, pp.~2825--2830, 2011.

\end{thebibliography}

% ─────────────────────────────────────────────────────────────────────────────
\label{sec:ai-disclaimer}

Parts of this project have been developed with the assistance of OpenAI's ChatGPT (GPT-5). The AI was used to support the development of project ideas, the structuring of methodological workflows, the drafting of descriptive texts, and the identification of relevant datasets and references. All content produced with AI assistance has been carefully reviewed, edited, and validated by me. I take full responsibility for the final content and its accuracy, relevance, and academic integrity.

\subsection*{Using AI (for students)}
Generative AI tools (such as ChatGPT, Claude, Mistral, or similar models) may be used in this project, both as an object of investigation and as a tool to support the development process. Students are encouraged to explore how these models function, interact with them creatively, and leverage them as inspiration or assistance in ideation, drafting, or experimentation.

However, AI should not be used as a substitute for original work. The responsibility for the structure, reasoning, and understanding of the project remains entirely with the student.

If generative AI has been used at any stage of the project, it is mandatory to include a disclaimer clearly specifying:
\begin{itemize}
    \item Which models have been used (e.g., GPT-4, Claude 3, etc.)
    \item For what purposes (e.g., drafting text, summarizing ideas, generating code or examples)
    \item To what extent the outputs were modified, verified, or integrated into the final submission
\end{itemize}

The project will be assessed not only based on its output, but also on the student's ability to explain and justify all choices made. A final interview will evaluate the depth of understanding, and any lack of clarity or over-reliance on AI-generated material without proper insight may negatively affect the evaluation.

\end{document}
